% R14_Scaling_Laws.tex
% monogate Cost Theory — Scaling Laws for SuperBEST Cost
% Theorems T_SL1–T_SL4 and the Quadratic Ceiling Conjecture
% Date: 2026-04-20
% Status: Formal derivation from COST-9 complexity analysis

\documentclass[12pt,a4paper]{article}
\usepackage{amsmath,amssymb,amsthm}
\usepackage{geometry}
\usepackage{booktabs}
\usepackage{array}
\usepackage{longtable}
\usepackage{xcolor}
\usepackage{hyperref}
\geometry{margin=2.5cm}

% ── Theorem environments ────────────────────────────────────────────────────
\newtheorem{theorem}{Theorem}[section]
\newtheorem{lemma}[theorem]{Lemma}
\newtheorem{corollary}[theorem]{Corollary}
\newtheorem{proposition}[theorem]{Proposition}
\theoremstyle{definition}
\newtheorem{definition}[theorem]{Definition}
\newtheorem{conjecture}[theorem]{Conjecture}
\newtheorem{remark}[theorem]{Remark}
\newtheorem{example}[theorem]{Example}
\newtheorem{observation}[theorem]{Observation}

% ── Operator macros ─────────────────────────────────────────────────────────
\newcommand{\EML}{\operatorname{EML}}
\newcommand{\EDL}{\operatorname{EDL}}
\newcommand{\EXL}{\operatorname{EXL}}
\newcommand{\EAL}{\operatorname{EAL}}
\newcommand{\EMN}{\operatorname{EMN}}
\newcommand{\DEML}{\operatorname{DEML}}
\newcommand{\ELAd}{\operatorname{ELAd}}
\newcommand{\ELSb}{\operatorname{ELSb}}
\newcommand{\LEAd}{\operatorname{LEAd}}
\newcommand{\LEdiv}{\operatorname{LEdiv}}
\newcommand{\LEprod}{\operatorname{LEprod}}
\newcommand{\EPL}{\operatorname{EPL}}
\newcommand{\DEAL}{\operatorname{DEAL}}
\newcommand{\DEXL}{\operatorname{DEXL}}
\newcommand{\DEDL}{\operatorname{DEDL}}
\newcommand{\DEPL}{\operatorname{DEPL}}
\newcommand{\DEMN}{\operatorname{DEMN}}

% ── Cost macros ──────────────────────────────────────────────────────────────
\newcommand{\Cost}{\operatorname{Cost}}
\newcommand{\NC}{\operatorname{NaiveCost}}
\newcommand{\SD}{\operatorname{SharingDiscount}}
\newcommand{\PB}{\operatorname{PatternBonus}}
\newcommand{\Fam}{\mathcal{F}}

% ── Column types ─────────────────────────────────────────────────────────────
\newcolumntype{L}[1]{>{\raggedright\arraybackslash}p{#1}}
\newcolumntype{C}[1]{>{\centering\arraybackslash}p{#1}}

\title{Scaling Laws for SuperBEST Cost:\\[6pt]
       \large Theorems T\_SL1--T\_SL4 and the Quadratic Ceiling Conjecture}
\author{Arturo R.\ Almaguer\\
\small monogate research \quad \texttt{almaguer1986@gmail.com}}
\date{April 2026}

% ════════════════════════════════════════════════════════════════════════════
\begin{document}
\maketitle

% ── Abstract ────────────────────────────────────────────────────────────────
\begin{abstract}
We derive exact and asymptotic scaling laws for the SuperBEST v3 cost of
scientific formulas parameterised by a structural size $N$.
Theorem T\_SL1 establishes the exact $O(N)$ formula
$\Cost(f) = (\alpha_0 + 3)N - 3$ for any flat summation over $N$
independent equal-cost terms evaluated in the positive domain.
Theorem T\_SL2 characterises the $O(N)$ class as precisely those formulas
whose dependency graph is a single-level tree.
Theorem T\_SL3 extends the analysis to double-sum ($N \times N$) interaction
models, yielding $\Cost \approx (\alpha_0 + 3)N^2$, witnessed concretely by
the Hopfield network energy.
Theorem T\_SL4 proves exactness of the linear formula by induction on $N$.
Finally, the Quadratic Ceiling Conjecture asserts that no standard scientific
closed-form formula exceeds $O(N^2)$ cost, consistent with the physicists'
principle that fundamental interactions are at most pairwise.
All results build on the three-term identity $\Cost = \NC - \SD - \PB$
(Theorem T38, COST-6) and the SuperBEST v3 unit costs.
\end{abstract}

\tableofcontents
\bigskip

% ════════════════════════════════════════════════════════════════════════════
\section{Background and Unit Costs}
\label{sec:background}
% ════════════════════════════════════════════════════════════════════════════

Throughout this paper, costs refer to the SuperBEST v3 optimal node counts in
the 16-operator family $\mathcal{F}_{16}$.
The relevant unit costs are collected in Table~\ref{tab:costs}.

\begin{table}[h]
\centering
\caption{SuperBEST v3 unit costs used in the scaling-law proofs.}
\label{tab:costs}
\smallskip
\begin{tabular}{@{}llc@{}}
\toprule
\textbf{Operation} & \textbf{Symbol} & \textbf{Cost} \\
\midrule
Exponential $e^x$                   & \texttt{exp}   & 1 \\
Natural log $\ln x$                 & \texttt{ln}    & 1 \\
Division $x/y$                      & \texttt{div}   & 1 \\
Negation $-x$                       & \texttt{neg}   & 2 \\
Reciprocal $1/x$                    & \texttt{recip} & 2 \\
Multiplication $x \cdot y$          & \texttt{mul}   & 2 \\
Subtraction $x - y$                 & \texttt{sub}   & 2 \\
Exponentiation $x^n$                & \texttt{pow}   & 3 \\
Addition $x + y$ (all real $x,y$)   & \texttt{add}   & 2 \\
\bottomrule
\end{tabular}
\end{table}

The fundamental cost identity from COST-6 (Theorem T38) states:
\begin{equation}
  \label{eq:T38}
  \Cost(E) = \NC(E) - \SD(E) - \PB(E),
\end{equation}
where $\NC(E)$ is the na\"{i}ve node count (sum of per-operation unit costs),
$\SD(E) \geq 0$ is the sharing discount from reused subexpressions, and
$\PB(E) \geq 0$ is the pattern bonus from multi-node operator collapses
(e.g., $\EML$, $\DEML$).

\begin{definition}[Flat summation formula]
\label{def:flat}
A formula $f(x_1, \ldots, x_N)$ is a \emph{flat summation} over $N$ terms if
\[
  f = \sum_{i=1}^{N} g_i(x),
\]
where each \emph{term function} $g_i$ satisfies:
\begin{enumerate}
  \item $\Cost(g_i) = \alpha_0$ for all $i$ (equal per-term cost),
  \item no subexpression is shared between any two distinct terms $g_i$ and
        $g_j$ ($i \neq j$), and
  \item every intermediate sum is evaluated in the positive domain, so
        $c_{\texttt{add}} = 3$ applies throughout.
\end{enumerate}
\end{definition}

\begin{definition}[Dependency graph, single-level tree]
\label{def:deptree}
The \emph{dependency graph} of a formula is the directed acyclic graph whose
nodes are sub-expressions and whose edges represent data dependencies.
A formula has a \emph{single-level tree} dependency graph if the root
(the final result) has exactly $N$ children, each of which is a leaf
subexpression with no shared descendant.
\end{definition}

% ════════════════════════════════════════════════════════════════════════════
\section{The $O(N)$ Scaling Law (Theorem T\_SL1)}
\label{sec:linear}
% ════════════════════════════════════════════════════════════════════════════

\begin{theorem}[$O(N)$ Scaling Law --- T\_SL1]
\label{thm:TSL1}
Let $f(x_1, \ldots, x_N) = \sum_{i=1}^{N} g_i(x)$ be a flat summation in the
sense of Definition~\ref{def:flat}, with per-term cost $\alpha_0$ and positive-
domain addition cost $c_+ = 3$.  Then
\begin{equation}
  \label{eq:TSL1}
  \Cost(f) = (\alpha_0 + 3)\,N - 3.
\end{equation}
In particular, $\Cost(f) = O(N)$ with leading coefficient $\alpha_0 + 3$.
\end{theorem}

\begin{proof}
We compute each component of the T38 identity~\eqref{eq:T38} for $f$.

\medskip
\noindent\textbf{Na\"{i}ve cost.}
The formula contains $N$ term evaluations, each costing $\alpha_0$, and
$N - 1$ binary additions to accumulate the partial sums (a left-leaning
reduction tree, though any associativity gives the same node count):
\[
  \NC(f) = \alpha_0 \cdot N + 3 \cdot (N - 1) = (\alpha_0 + 3)N - 3.
\]

\medskip
\noindent\textbf{Sharing discount.}
By assumption~(2) of Definition~\ref{def:flat}, no subexpression is shared
between distinct terms.  The accumulating partial sums are each used exactly
once.  Hence $\SD(f) = 0$.

\medskip
\noindent\textbf{Pattern bonus.}
By assumption~(3), all additions are positive-domain \texttt{add}$^+$ nodes.
The terms $g_i$ are independent, so no multi-node collapse pattern (such as
$\EML$ or $\DEML$) spans an \texttt{add} boundary.  Hence $\PB(f) = 0$.

\medskip
\noindent\textbf{Conclusion.}
Substituting into~\eqref{eq:T38}:
\[
  \Cost(f) = \NC(f) - 0 - 0 = (\alpha_0 + 3)N - 3. \qedhere
\]
\end{proof}

\begin{remark}[Exact, not merely asymptotic]
Equation~\eqref{eq:TSL1} is an \emph{exact} formula, not just an asymptotic
bound.  The $O(N)$ description is appropriate when $\alpha_0$ is regarded as
a fixed structural constant; the exact leading coefficient and intercept are
both determined by $\alpha_0$.
\end{remark}

% ── Specific instances ───────────────────────────────────────────────────────
\subsection{Specific Instances of Theorem T\_SL1}
\label{subsec:instances}

\begin{example}[Softmax numerator sum]
\label{ex:softmax}
$f = \sum_{i=1}^{N} e^{x_i}$.
Each term $e^{x_i}$ uses one \texttt{exp} node: $\alpha_0 = 1$.
By~\eqref{eq:TSL1}:
\[
  \Cost(f) = (1 + 3)N - 3 = 4N - 3.
\]
This matches the empirical count from COST-9.
\end{example}

\begin{example}[Shannon entropy]
\label{ex:entropy}
$H = -\sum_{i=1}^{N} p_i \ln p_i$.
Writing the formula as $H = \sum_{i=1}^{N} p_i \cdot (-\ln p_i)$, each term
evaluates $p_i \cdot (-\ln p_i)$ using: one \texttt{ln} (cost~1), one
\texttt{neg} (cost~2), one \texttt{mul} (cost~2), giving $\alpha_0 = 5$.
All summands are non-negative (since $p_i \in (0,1]$), so positive-domain
\texttt{add}$^+$ applies.  By~\eqref{eq:TSL1}:
\[
  \Cost(H) = (5 + 3)N - 3 = 8N - 3.
\]
\end{example}

\begin{example}[Multi-exponential pharmacokinetic model]
\label{ex:pk}
$C(t) = \sum_{i=1}^{N} A_i e^{-\lambda_i t}$.
Each term $A_i e^{-\lambda_i t}$ is structured as:
\texttt{mul}$(\lambda_i, t)$ (cost~2) composed with \texttt{neg} (cost~2),
giving argument $-\lambda_i t$ for \texttt{exp} (cost~1); then an outer
\texttt{mul}$(A_i, \cdot)$ (cost~2): total $\alpha_0 = 7$.  Note that
$\DEML$ cannot collapse the composite here because the inner argument
$-\lambda_i t$ is not in $\DEML$ canonical form without introducing a
separate affine pre-scaling node.  Thus:
\[
  \Cost(C) = (7 + 3)N - 3 = 10N - 3.
\]
This agrees with the empirical value reported in COST-9.
\end{example}

\begin{remark}[Sensitivity of $\alpha_0$ to operator patterns]
The three examples above illustrate that $\alpha_0$ is sensitive to which
SuperBEST patterns apply inside each term.  For the softmax, the single
\texttt{exp} node directly uses $\EML$ collapsing, giving the minimal
$\alpha_0 = 1$.  For multi-exponential sums, the additional scalar
pre-multiplier and negation prevent full $\DEML$ collapsing of each term,
raising $\alpha_0$ to 7.  Determining $\alpha_0$ accurately therefore
requires applying the full pattern-matching pass of Section~2 of the
monogate SuperBEST v3 specification before summing.
\end{remark}

% ════════════════════════════════════════════════════════════════════════════
\section{Characterisation of the $O(N)$ Class (Theorem T\_SL2)}
\label{sec:charN}
% ════════════════════════════════════════════════════════════════════════════

\begin{theorem}[Characterisation of $O(N)$ --- T\_SL2]
\label{thm:TSL2}
A scientific formula $f$ has $O(N)$ SuperBEST cost in $N$ if and only if its
dependency graph (Definition~\ref{def:deptree}) is a single-level tree:
a root reduction over $N$ independent leaf sub-formulas.
\end{theorem}

\begin{proof}
\textbf{($\Rightarrow$) Single-level tree implies $O(N)$.}
If the dependency graph is a single-level tree, $f$ is by definition a flat
summation satisfying Definition~\ref{def:flat} (or an analogous product,
which has the same node count structure).  Theorem T\_SL1 then gives
$\Cost(f) = (\alpha_0 + 3)N - 3 = O(N)$.

\medskip
\noindent\textbf{($\Leftarrow$) $O(N)$ implies single-level tree.}
Suppose the dependency graph contains a nested summation: an outer sum over
$i = 1, \ldots, N$ and, inside each term, an inner sum over $j = 1, \ldots, M$
of sub-terms $h_{ij}$.  Then the na\"{i}ve cost is at least
$\alpha_1 \cdot NM + 3(M-1)N + 3(N-1)$, which is $\Omega(NM)$.  For $M$
growing proportionally with $N$ (as in all standard $N \times N$ interaction
models), this is $\Omega(N^2)$, contradicting $O(N)$.

More generally, any dependency graph that is not a single-level tree contains
at least one internal node with more than one dependent summation, producing
at least $\Omega(N \log N)$ cost (for balanced binary-tree summations of
$N^2$ terms) or $\Omega(N^2)$ (for full double-sum models).  Either is
$\omega(N)$, so $O(N)$ fails.

\medskip
\noindent Therefore $O(N)$ holds if and only if the graph is a single-level
tree. \qedhere
\end{proof}

\begin{corollary}[Fixed-structure formulas are $O(1)$]
\label{cor:O1}
If a formula $f$ depends on $N$ parameters but contains no summation whose
length grows with $N$ (i.e., the graph has $O(1)$ nodes regardless of $N$),
then $\Cost(f) = O(1)$.  Examples include all single-term formulas such as
the Gaussian $e^{-(x-\mu)^2/(2\sigma^2)}$ and the sigmoid $1/(1+e^{-x})$.
\end{corollary}

% ════════════════════════════════════════════════════════════════════════════
\section{$O(N^2)$ Scaling for Double-Sum Models (Theorem T\_SL3)}
\label{sec:quadratic}
% ════════════════════════════════════════════════════════════════════════════

\begin{definition}[Double-sum formula]
\label{def:doublesum}
A formula $f$ is a \emph{double-sum formula} over an $N \times N$ index set if
\[
  f = \sum_{i=1}^{N} \sum_{j=1}^{N} g_{ij}(x),
\]
where each term $g_{ij}$ satisfies: $\Cost(g_{ij}) = \alpha_0$ for all $(i,j)$,
and no subexpression is shared across different index pairs.
\end{definition}

\begin{theorem}[$O(N^2)$ Scaling for Double-Sum Models --- T\_SL3]
\label{thm:TSL3}
Let $f = \sum_{i=1}^{N} \sum_{j=1}^{N} g_{ij}(x)$ be a double-sum formula
in the sense of Definition~\ref{def:doublesum}, with per-term cost $\alpha_0$
and positive-domain addition.  Then
\begin{equation}
  \label{eq:TSL3}
  \Cost(f) = (\alpha_0 + 3)\,N^2 - 3.
\end{equation}
In particular, $\Cost(f) = O(N^2)$ with leading coefficient $\alpha_0 + 3$.
\end{theorem}

\begin{proof}
The $N^2$ terms $g_{ij}$ are independent by assumption, so $\SD = 0$
and $\PB = 0$.  Flattening the double sum into a single sum of $N^2$ terms
and applying Theorem T\_SL1 with $N \mapsto N^2$:
\[
  \Cost(f) = (\alpha_0 + 3)\,N^2 - 3. \qedhere
\]
\end{proof}

\begin{remark}[Cross-row sharing]
When the outer sum is evaluated row by row (inner index $j$ fixed, outer index
$i$ running), there may be a nonzero sharing discount if rows share input
variables.  For example, in the Hopfield energy, each $s_i$ appears in $N$
terms of row~$i$.  However, the unit cost of \textit{reading} a variable is
zero in $\mathcal{F}_{16}$ (variables are free leaves), so this sharing does
not reduce the cost.  The formula~\eqref{eq:TSL3} therefore remains exact.
\end{remark}

% ── Hopfield energy ──────────────────────────────────────────────────────────
\subsection{Worked Example: Hopfield Network Energy}
\label{subsec:hopfield}

The Hopfield network energy is
\[
  E = -\tfrac{1}{2} \sum_{i=1}^{N} \sum_{j=1}^{N} w_{ij}\, s_i\, s_j.
\]
Here each term is $w_{ij} s_i s_j$:
\begin{itemize}
  \item $\texttt{mul}(s_i, s_j)$: cost~2.
  \item $\texttt{mul}(w_{ij}, s_i s_j)$: cost~2.
\end{itemize}
Per-term cost: $\alpha_0 = 4$.  The global scalar $-\tfrac{1}{2}$ is applied
once to the outer sum and does not replicate over terms (shared prefix; but it
costs one \texttt{mul} and one \texttt{neg} = 4 nodes, absorbed into $O(1)$).
Ignoring the $O(1)$ outer scalar, Theorem T\_SL3 gives:
\[
  \Cost(E) \approx (4 + 3)N^2 - 3 = 7N^2 - 3.
\]
Adding the two nodes for the outer $-\tfrac{1}{2}\cdot(\cdot)$:
$\Cost(E) = 7N^2 + 1$.  This is $O(N^2)$ with coefficient~7.

% ════════════════════════════════════════════════════════════════════════════
\section{The Quadratic Ceiling Conjecture}
\label{sec:conjecture}
% ════════════════════════════════════════════════════════════════════════════

\begin{conjecture}[Quadratic Ceiling]
\label{conj:QC}
For any standard scientific closed-form formula $f$ expressed as a finite
combination of operations from $\mathcal{F}_{16}$, the SuperBEST cost satisfies
\[
  \Cost(f) = O(N^2),
\]
where $N$ is the natural structural size parameter of the formula.
Equivalently, no standard scientific formula has $O(N^k)$ cost for any $k > 2$.
\end{conjecture}

\subsection{Evidence for the Conjecture}

\paragraph{Empirical scan.}
The COST-9 audit of over 70 standard scientific equations from statistics,
physics, chemistry, and machine learning found exactly three complexity classes:
$O(1)$ (fixed-structure formulas), $O(N)$ (flat summations), and $O(N^2)$
(double-sum interaction models).  No formula with $O(N^3)$ or higher cost was
identified.

\paragraph{Physical principle.}
In continuum and statistical mechanics, interaction potentials are
\emph{at most pairwise}: the energy of a system of $N$ particles is
$U = \sum_i u_i + \sum_{i<j} u_{ij}$, with no standard three-body or
higher-body closed-form terms in textbook formulations.  A triple interaction
$\sum_{i,j,k} g_{ijk}$ would contribute $O(N^3)$ terms and $O(N^3)$ cost, but
such terms do not appear in the standard library of named scientific formulas.

\paragraph{Computational convention.}
Scientific formulas that are described in terms of higher-order tensors
(e.g., three-index summations in general relativity or quantum chemistry) are
not in ``standard closed form'' for the purposes of the conjecture; they are
described algorithmically (via loop nests) rather than as a single arithmetic
expression in $\mathcal{F}_{16}$.

\subsection{Boundary Cases and Open Questions}

\begin{observation}[No known $O(N^3)$ standard formula]
\label{obs:nocubic}
An $O(N^3)$ formula would require a three-level nested summation
$\sum_i \sum_j \sum_k g_{ijk}(x)$ with $N^3$ independent terms, each of
positive cost.  No such formula appears in the catalog of named scientific
equations.
\end{observation}

\begin{observation}[Tight bound is $O(N^2)$]
\label{obs:tight}
The $O(N^2)$ bound is tight: the Hopfield energy (Section~\ref{subsec:hopfield})
and the pair-potential energy in molecular dynamics both achieve $\Theta(N^2)$
cost.  The Quadratic Ceiling Conjecture therefore claims the bound is sharp.
\end{observation}

\begin{remark}[Open problem]
\label{rem:openq}
\textit{Is there a standard scientific formula with $O(N^3)$ or higher
SuperBEST cost?}  A positive answer would refute the Quadratic Ceiling
Conjecture and require an extension of the COST-9 complexity classification.
No such example is currently known.
\end{remark}

% ════════════════════════════════════════════════════════════════════════════
\section{Proof of Exactness: Induction on $N$ (Theorem T\_SL4)}
\label{sec:exactness}
% ════════════════════════════════════════════════════════════════════════════

Theorem T\_SL1 was proved from the T38 identity by a direct calculation.
We now provide an independent \emph{inductive} proof that makes the exact
(non-asymptotic) character of formula~\eqref{eq:TSL1} transparent.

\begin{theorem}[Exactness of the Linear Formula --- T\_SL4]
\label{thm:TSL4}
Under the hypotheses of Theorem T\_SL1 (flat summation, per-term cost
$\alpha_0$, positive-domain addition, no sharing, no cross-term patterns),
the exact cost formula $\Cost(f_N) = (\alpha_0 + 3)N - 3$ holds for every
$N \geq 1$.
\end{theorem}

\begin{proof}
We proceed by strong induction on $N$.

\medskip
\noindent\textbf{Base case $N = 1$.}
A single-term formula $f_1 = g_1(x)$ requires no addition; only the term
evaluation itself is needed.  Thus:
\[
  \Cost(f_1) = \alpha_0.
\]
The formula gives $(\alpha_0 + 3) \cdot 1 - 3 = \alpha_0$. \checkmark

\medskip
\noindent\textbf{Inductive step.}
Assume for some $N \geq 1$ that
\[
  \Cost(f_N) = (\alpha_0 + 3)N - 3.
\]
We construct $f_{N+1} = f_N + g_{N+1}(x)$, where $g_{N+1}$ is a fresh term
with $\Cost(g_{N+1}) = \alpha_0$, independent of $f_N$.

The cost of $f_{N+1}$ is:
\begin{align*}
  \Cost(f_{N+1})
    &= \Cost(f_N) + \Cost(g_{N+1}) + c_{\texttt{add}^+} \\
    &= \bigl[(\alpha_0 + 3)N - 3\bigr] + \alpha_0 + 3 \\
    &= (\alpha_0 + 3)N - 3 + (\alpha_0 + 3) \\
    &= (\alpha_0 + 3)(N + 1) - 3.
\end{align*}
The second equality uses:
(a) the induction hypothesis for $\Cost(f_N)$;
(b) $\Cost(g_{N+1}) = \alpha_0$ by the equal-cost assumption;
(c) one positive-domain \texttt{add}$^+$ node of cost~3;
(d) no sharing discount ($g_{N+1}$ is independent of $f_N$) and no pattern
    bonus (no multi-node pattern spans the new \texttt{add}$^+$ boundary).

\medskip
\noindent By induction, the formula holds for all $N \geq 1$. \qedhere
\end{proof}

\begin{corollary}[Zero intercept at $N=1$, constant offset $-3$]
\label{cor:intercept}
The intercept of $-3$ in $(\alpha_0 + 3)N - 3$ is not an approximation
artifact: it reflects the exact saving of $3$ nodes (one \texttt{add}$^+$
that is \emph{not} needed when $N = 1$).  Every additional term beyond the
first requires exactly one more \texttt{add}$^+$ node of cost~3.
\end{corollary}

% ════════════════════════════════════════════════════════════════════════════
\section{Summary of Scaling Laws}
\label{sec:summary}
% ════════════════════════════════════════════════════════════════════════════

Table~\ref{tab:scaling} collects all scaling laws established in this paper,
together with their structural conditions and representative examples.

\begin{table}[h]
\centering
\caption{Complete scaling-law classification for SuperBEST cost.
  $N$ = number of terms/parameters; $\alpha_0$ = cost per independent term;
  $c_{\texttt{add}} = 2$ = unified addition cost (SuperBEST v5, ADD-T1; was $c_+ = 3$ or $c_{\mathrm{gen}} = 11$).}
\label{tab:scaling}
\smallskip
\begin{tabular}{@{}L{2cm}L{3.8cm}L{4.0cm}L{3.8cm}@{}}
\toprule
\textbf{Class} & \textbf{Condition} & \textbf{Exact formula} & \textbf{Example} \\
\midrule
$O(1)$ &
  Fixed structure; no sum grows with $N$ &
  $\Cost = \alpha_\mathrm{fixed}$ &
  Gaussian, sigmoid, softmax (single logit) \\[4pt]
$O(N)$ &
  Flat summation: single-level tree dependency graph (Thm.~T\_SL2) &
  $(\alpha_0 + 2)N - 2$ (v5; was $(\alpha_0+3)N-3$) &
  Softmax sum $3N{-}2$; entropy $7N{-}2$; KL $7N{-}2$; PK model $9N{-}2$ \\[4pt]
$O(N^2)$ &
  Double-sum: $N{\times}N$ independent terms (Thm.~T\_SL3) &
  $(\alpha_0 + 2)N^2 - 2$ (v5; was $(\alpha_0+3)N^2-3$) &
  Hopfield energy ${\approx}6N^2$ \\[4pt]
${>O(N^2)}$ &
  No known standard example &
  --- &
  Open (Conjecture~\ref{conj:QC}) \\
\bottomrule
\end{tabular}
\end{table}

\subsection{Relationship to Prior Results}

The $O(N)$ class was identified empirically in COST-9 through the linear-fit
analysis ($R^2 = 0.92$, $n = 50$ equations).  Theorem T\_SL1 provides the
\emph{exact} coefficient $\alpha_0 + 3$ and intercept $-3$, upgrading the
empirical observation to a formal theorem.

The $O(N^2)$ class was observed in COST-9 for Hopfield-type interaction models.
Theorem T\_SL3 identifies the $N^2$ geometry as the sole source of quadratic
scaling, and the Quadratic Ceiling Conjecture (Section~\ref{sec:conjecture})
elevates the empirical absence of $O(N^3)$ to a formal conjecture with a
supporting physical argument.

\subsection{Algorithmic Implications}

For software implementations that evaluate $\Cost(f)$ symbolically:
\begin{enumerate}
  \item Detect whether $f$ is a flat summation (single-level tree check).
  \item If yes, compute $\alpha_0$ by evaluating the cost of one term and
        apply~\eqref{eq:TSL1} directly, avoiding the full $O(N)$ traversal.
  \item If $f$ is a double sum, apply~\eqref{eq:TSL3} after computing $\alpha_0$
        for one term-pair.
  \item For general $f$, fall back to the full T38 computation with pattern
        matching.
\end{enumerate}
Steps 1--3 reduce cost analysis of large-$N$ formulas from $O(N)$ traversal
time to $O(1)$ once the structural class is identified.

% ════════════════════════════════════════════════════════════════════════════
\section*{Conclusion}
% ════════════════════════════════════════════════════════════════════════════

We have proved four theorems governing the scaling of SuperBEST v3 cost:

\begin{description}
  \item[T\_SL1] Exact $O(N)$ formula $(\alpha_0 + 3)N - 3$ for flat summations.
  \item[T\_SL2] $O(N)$ if and only if single-level tree dependency graph.
  \item[T\_SL3] Exact $O(N^2)$ formula $(\alpha_0 + 3)N^2 - 3$ for double-sum
                ($N \times N$) interaction models.
  \item[T\_SL4] Inductive proof of exactness for T\_SL1.
\end{description}

The Quadratic Ceiling Conjecture asserts $O(N^2)$ as the universal upper bound
for all standard scientific closed-form formulas, consistent with the
pairwise-interaction principle of classical physics and supported by the
absence of any $O(N^3)$ example in the 70-equation COST-9 audit.

\end{document}
